\section{开发计划与进度安排}
\begin{frame}{已有基础与后续工作}
\begin{block}{已记录的原型进展}
仓库已具备工作台、共享执行服务、浏览器操作与证据报告闭环。验收记录记载：一个真实 DeepSeek + 4gaBoards 创建项目场景通过，使用 8 次操作。
\end{block}
\begin{itemize}
\item 该结果来自单个自建场景，只用于说明原型可行性。
\item 教师的 5 个调试场景尚未取得，完整课程场景尚未评测。
\item 后续重点是跨应用适配、复杂交互恢复、标准答案核对与可复现实验。
\item 新成员安装流程仍需在另一台机器验证，协作平台也需完成配置。
\end{itemize}
\vfill
\footnotesize 来源：仓库 README 与 docs/acceptance.md；不将离线模拟结果计作真实模型成绩。
\end{frame}

\begin{frame}{拟定进度与阶段交付}
\small
\begin{tabularx}{\textwidth}{lXX}
\toprule
阶段 & 主要工作 & 里程碑与验收产物 \\
\midrule
第 1 周 & 确认课程约束、成员与场景，建立远端协作，复核已有原型 & 需求清单、责任分配、环境初始化记录。\\
第 2--3 周 & 完善定位、结构化动作、预算控制与证据审查 & 代表场景可运行，典型异常可解释且有证据。\\
第 4 周 & 完成工作台联调和三个应用的测试环境适配 & 界面演示流程、环境版本清单与联调记录。\\
第 5 周 & 执行真实场景评测及对照实验，分析失败原因 & 逐场景结果、标准答案对照及策略比较。\\
第 6 周 & 修复问题、验证部署复现、完善报告并排练答辩 & 源码、部署文档、实验报告及演示材料。\\
\bottomrule
\end{tabularx}
\vfill
\footnotesize 暂按开题后六周安排；实际日期以课程节点为准。各阶段由对应责任人推进，每周同步进度与阻塞项。
\end{frame}
